\chapter*{Preface}
Volume~6 ended with the prime number theorem, its error term, and the first sieves. It left the primes looking regular from far away and mysterious from close up. Is there always a prime between $x$ and $x+\sqrt x$? How evenly are the primes spread among the residue classes of a modulus that grows with $x$? How many ways are there to write a large number as a sum of $k$-th powers, or an odd number as a sum of three primes? Are there infinitely many pairs of primes at bounded distance? Does the M\"obius function cancel on average for reasons that have nothing to do with the zeros of $\zeta$? This book proves the answers that are known, most of them found in the last hundred years and several in the last fifteen.

The first part, \emph{Zeros and Sieves}, works with the zeros of $\zeta$ and of Dirichlet $L$-functions, and with sieves. Few zeros of $\zeta$ lie far from the critical line, and this is enough to find primes in every interval $\left]x,x+x^{0.76}\right]$. Siegel's theorem shows that an $L$-function cannot vanish too close to $1$, at the price of a constant that nobody can compute, and it gives the growth of class numbers promised in Volume~13. Sieves meet the parity problem, which explains why they cannot find primes alone. The large sieve proves the Bombieri--Vinogradov theorem: on average over moduli up to $\sqrt x$, primes in progressions behave as the generalised Riemann hypothesis predicts. Vinogradov's method bounds exponential sums over primes.

The second part, \emph{Additive Problems and Bounded Gaps}, puts these tools to work. The circle method of Hardy and Littlewood counts sums of $k$-th powers, and the book proves that every large number is a sum of $\max(4k,2^k+1)$ of them. It counts sums of three primes, and gives Vinogradov's theorem. Then the sieve weights of Goldston, Pintz and Y\i ld\i r\i m, in Maynard's form, prove that the gaps between consecutive primes are bounded infinitely often. The bound proved on the page, with every constant written out and checked, is $4580$; the records $600$ and $246$ are quoted.

The third part, \emph{Structure in the Primes}, asks what the primes and the multiplicative functions look like beyond their averages. Hal\'asz's theorem says when a multiplicative function has a mean value. Additive combinatorics brings the Pl\"unnecke--Ruzsa inequality, the Balog--Szemer\'edi--Gowers theorem and the sum-product phenomenon; Roth's theorem finds three-term progressions in every dense set, and Behrend shows how sparse such a set can be. Gowers norms count longer progressions and explain the architecture of the Green--Tao theorem. The book ends on the critical line, with Hardy's theorem that infinitely many zeros lie on it, the mean square of $\zeta$ there, and Montgomery's pair correlation of the zeros, where number theory meets random matrices.

The aim throughout is to understand why the theorems are true. Every result is proved, except deep theorems that are quoted, each marked as such where it appears and never used silently. The list is not short: it includes the zero-free region of Vinogradov and Korobov, Chen's theorem, the Vinogradov mean value theorem, Weil's bound for Kloosterman sums, Helfgott's theorem on three primes, Zhang's theorem, the Matom\"aki--Radziwi\l\l{} theorem, Freiman's theorem in $\mathbb Z$, the theorems of Szemer\'edi and of Kelley and Meka, the higher Gowers inverse theorem, the Green--Tao theorem, and a positive proportion of zeros on the critical line. Results that need a hypothesis, such as the generalised Riemann hypothesis or a level of distribution beyond $\frac12$, say so in their headings. The Siegel--Walfisz theorem, the Bombieri--Vinogradov theorem and the three primes theorem share one ineffective constant, coming from Siegel's theorem. Pause and think boxes ask you to guess before reading on. Each chapter ends with exercises, graded from warm-ups ($\star$) to small gems ($\star\star\star$), each with a hint and a complete solution that first explains how to find the idea and then writes it up.

The book assumes Volume~6, and uses Volumes~5, 8, 11, 12 and 13 where they help. Its beauty is in the way the primes, which no formula describes, obey laws that can be proved: deterministic, yet statistically as regular as a random sequence. Each chapter tries to show how much of that regularity can be reached, and with which idea. Volume~15 takes the question further, and asks which of these regularities a random model predicts, which a theorem guarantees, and which are still only believed.
